\section{Paths for a complete native cell attached under an old algebra}
\label{sec:136-relative-cell-paths}

\paragraph{Archive and exact scope.}
The relevant uniform conditions are cote136, physical pp.~117 and119,
footers116 and118, scans117 and119; the common-source question is
physical p.~121, footer120, scan121. This result treats an actual
boundary-constrained pushout under an existing algebra. Its hypothesis
includes the complete coherent boundary data in that algebra. It does
not construct those data from a native input alone. The proof uses
ARCH-136-DERIVED-TAIL-COMPARISON and ARCH-136-TAIL-EXT-VANISHING,
and the full strict presentation of
ARCH-136-COHERENT-PD-REALIZATION. It is a relative version of the
prescribed-path argument, with the whole boundary in place of a graph.

Work over \(R=\mathbb Z\) or \(\mathbb Z_{(\ell)}\), and put
\(S=\Gamma_R(T)\). Let \(q\geq1\) be odd, \(p\geq2\) an
integer, \(m>0\), and
\[
             D=p(q+1),\qquad v=pm,\qquad
             X=\Delta[n],\quad B=\partial\Delta[n],
             \quad n\geq2.                              \tag{136.RC1}
\]
Fix \(M\geq1\) and a free weight-graded resolution
\(P_\bullet\xrightarrow{\epsilon}I_M\), all of whose
homogeneous generators have weights at least \(M\).

\paragraph{The actual relative presentation.}
For a simplicial set \(K\), define \(\mathcal N(K)\) as follows.
For every basis simplex \(\sigma\in K_r\), take an odd native
unit disk \(a_\sigma,b_\sigma\) of weights \(m,m\), form
degrees \(q,q+1\), and differential
\(da_\sigma=b_\sigma\), \(db_\sigma=0\).
Put
\[
 z_\sigma=\gamma_p(b_\sigma),\qquad
 e_\sigma=a_\sigma\gamma_{p-1}(b_\sigma),\qquad
                         de_\sigma=pz_\sigma.           \tag{136.RC2}
\]
Adjoin linear parameter copies
\[
 U_i:P_i[v]\longrightarrow\mathcal N(K)_r^{*,D-1-i}
       \ (0\leq i\leq D-1),\qquad
 V_i:P_i[v]\longrightarrow\mathcal N(K)_r^{*,D-2-i}
       \ (0\leq i\leq D-2)
\]
for that same basis simplex, with the following relations (the
simplex subscript is suppressed):
\[
\begin{aligned}
 dU_0&=\epsilon z,&
 dU_i&=-U_{i-1}\partial_i &&(1\leq i\leq D-1),\\
 dV_0&=pU_0-\epsilon e,&
 dV_i&=pU_i+V_{i-1}\partial_i &&(1\leq i\leq D-2),\\
 U_{D-1}\partial_D&=0,&
 pU_{D-1}+V_{D-2}\partial_{D-1}&=0.
\end{aligned}                                             \tag{136.RC3}
\]
Here \(\epsilon z\) means \(u\mapsto\epsilon(u)z\),
and similarly for \(\epsilon e\). Take the free graded strict
PD algebra on these generators, with odd generators exterior,
and quotient by the full PD ideal of the two bottom relations.
Its differential is the special differential satisfying
\(d\gamma_s(x)=\gamma_{s-1}(x)dx\) for even \(x\).
Equation (136.RC2), \(\epsilon\partial_1=0\), and
\(\partial^2=0\) give \(d^2=0\) and stability of this
ideal. For example the differential of the second bottom relation
is \(-pU_{D-2}\partial_{D-1}+pU_{D-2}\partial_{D-1}=0\);
the first is also closed. The special differential rule makes
the full generated PD ideal differential-stable.

Simplicial operators send the entire presentation for \(\sigma\)
to the presentation for its image simplex. Since a basis simplex
is sent to a basis simplex, the native divided-power relation is
respected without requiring additivity of \(\gamma_p\).
Thus \(\mathcal N(K)\) is an actual simplicial strict PD
differential algebra. Its maps to a target are exactly a native
section on \(K\), together with all the linear parameter sections
satisfying (136.RC3).

Let \(A\) be any such simplicial algebra, and prescribe a strict
map \(g_B:\mathcal N(B)\to A\). Equivalently, the native
boundary section \(a_B\in A^{m,q}(B)\) and all its coherent
parameters already belong to \(A(B)\) and satisfy (136.RC3).
Form the categorical pushout
\[
                E=A\amalg_{\mathcal N(B)}\mathcal N(X).
                                                               \tag{136.RC4}
\]
This means identifying every boundary generator with its prescribed
old value, not replacing the old algebra by a quotient chosen after
the fact. In each simplicial degree \(r\), the inclusion
\(B_r\subset X_r\) is an inclusion of bases, so (136.RC4)
is simply the strict PD coproduct over \(A_r\) of the independent
native coherent blocks indexed by \(X_r\setminus B_r\).
In particular it imposes no new relation inside \(A_r\).
Augmenting just those new blocks gives a differential graded
retraction in that degree; these retractions are not asserted to
commute with simplicial operators. This pointwise description
also verifies existence of the displayed pushout directly from
the presentation.

\paragraph{ARCH-136-RELATIVE-CELL-PATHS.}
Let \(C\) be admissible. Suppose its original input row
\(C_*^{m,q}\) is horizontally acyclic, and suppose
\[
                         M+v\geq W_C(D),                 \tag{136.RC5}
\]
where the comparison bound dominates the target's vertical and
horizontal cutoffs through form degree \(D+1\). Let
\(f_0,f_1:E\to C\) be two strict maps with the same
restriction \(f_A:A\to C\).

For every choice of an interpolation of the native input between
these maps, constant on the old full boundary, there is a single
obstruction class in
\[
 H^{D-1}\!\left(\Delta[n]\times\Delta[1],\,
  (\partial\Delta[n]\times\Delta[1])\cup
  (\Delta[n]\times\partial\Delta[1]);\,S^v/pS^v\right)
 \cong
 \begin{cases}
 S^v/pS^v,&n=D-2,\\
 0,&n\ne D-2.
 \end{cases}                                               \tag{136.RC6}
\]
It vanishes if and only if that input interpolation extends to
all coherent parameters with the prescribed endpoints and old
boundary values. In particular, when \(n\ne D-2\), there is
an actual strict simplicial path
\[
               H:E\longrightarrow P_{\Delta[1]}C,
       \qquad e_iH=f_i,\qquad
               H|_A=\operatorname{const}(f_A).             \tag{136.RC7}
\]
There is no additional lower bound on the original native weight
\(v\). Finitely many targets may be treated simultaneously,
using their maximum in (136.RC5); the construction of the source
pushout is unchanged.

\begin{proof}
Write \(I=\Delta[1]\), \(Z=X\times I\), and
\[
                         Y=(B\times I)\cup(X\times\partial I).
                                                               \tag{136.RC8}
\]
On \(Y\), the two endpoint native sections and the section
\(f_A(a_B)\) pulled back constantly along \(I\) agree on
overlaps. Acyclicity of \(C_*^{m,q}\) extends this prescribed
section to \(\widetilde a\in C^{m,q}(Z)\). Explicitly,
normalizing the free abelian groups on the inclusion \(Y\subset Z\)
gives degreewise free relative chains. Extend the associated chain
map into the input row one relative generator at a time: the
assigned image of its boundary is a cycle, and acyclicity supplies
a preimage. This proves existence for the entire simplicial pair,
rather than for unrelated simplexwise assignments. Fix any such
\(\widetilde a\) for the rest of the proof.

Use the total complex and sign convention
\[
\begin{split}
 \mathcal H(C(Z))^r
 &=\bigoplus_{i=0}^r\operatorname{Hom}^0_{S,\mathrm{gr}}
                        (P_i[v],C^{*,r-i}(Z)),\\
 (\delta f)_i&=df_i-(-1)^r f_{i-1}\partial_i.
\end{split}                                                \tag{136.RC9}
\]
Let \(\mathcal R\) be the restriction kernel to \(Y\).
It is exactly \(\mathcal H(C(Z,Y))\). All the restriction
maps needed in total degrees at most \(D\) are onto: every
\(P_i\) is free, every generator weight is at least \(M+v\),
and (136.RC5) puts the required form rows in the acyclic range.
The preceding relative-row extension argument applies separately
to each free coordinate. Arbitrary products of these surjections
are surjective; no lifting property of a nonfree final syzygy
has been used.

Put \(\widetilde b=d\widetilde a\),
\(\widetilde z=\gamma_p(\widetilde b)\), and
\(\widetilde e=\widetilde a\gamma_{p-1}(\widetilde b)\).
They give a closed element
\[
 \omega=(f_{\widetilde e},f_{\widetilde z})
       \in P_p(\mathcal H(C(Z)))^{D-1},\qquad
 d(x,y)=(\delta x-py,-\delta y).                           \tag{136.RC10}
\]
The already prescribed endpoint parameters and old boundary
parameters define \(\beta=(-V,-U)\) on \(Y\), of total
cone degree \(D-2\), satisfying \(d\beta=\omega|_Y\).
This uses all of (136.RC3), including the bottom equations.
They agree on intersections because the two maps agree on \(A\).
Lift \(\beta\) coordinatewise to \(\widetilde\beta\) on
\(Z\). Then
\[
          \kappa(\widetilde a)
             =[\omega-d\widetilde\beta]
                    \in H^{D-1}(P_p(\mathcal R))          \tag{136.RC11}
\]
is well-defined independently of that lift: two lifts differ by
a relative cochain. Its vanishing is precisely the condition that
one can add a relative cone cochain to \(\widetilde\beta\)
to obtain a primitive of \(\omega\), without changing its
restriction to \(Y\). Thus it is necessary and sufficient for
the asserted extension of the fixed native interpolation.

The derived-tail comparison identifies the cohomology of
\(\mathcal R\), through degree \(D\), with ordinary relative
cohomology of \((Z,Y)\) with coefficient complex
\[
 K=\operatorname{Hom}^0_{S,\mathrm{gr}}(P_\bullet[v],S).
\]
ARCH-136-TAIL-EXT-VANISHING supplies the canonical
quasi-isomorphism \(S^v[0]\to K\). Its terms are products
of torsion-free coefficient modules, so its Moore cone is
quasi-isomorphic to \(S^v/pS^v\) in degree zero. Applying
the cone exact sequence, or the natural comparison of cones,
therefore gives
\[
               H^{D-1}(P_p(\mathcal R))
                  \cong H^{D-1}(Z,Y;S^v/pS^v).           \tag{136.RC12}
\]
Only adjacent degrees \(D-1,D\) of the comparison are used.
All totalizations have finite diagonals in this range.

The relative normalized chain complex of (136.RC8) is chain
homotopy equivalent to the tensor product of the relative chains
of \((\Delta[n],\partial\Delta[n])\) and
\((\Delta[1],\partial\Delta[1])\). Each has one free
homology group, in degrees \(n\) and one respectively,
with split free contractions onto it. Consequently their relative
product has coefficient cohomology only in degree \(n+1\),
where it is the coefficient module itself. This proves (136.RC6).

When the obstruction vanishes, choose its relative primitive
\(\rho\), so that \(d\rho=\omega-d\widetilde\beta\),
and use \(\widetilde\beta+\rho\). Its components give
all the required \(U_i,V_i\) on \(Z\). The components
in the next resolution index, where the ordinary form degree
would be negative, are exactly
\(U_{D-1}\partial_D=0\) and
\(pU_{D-1}+V_{D-2}\partial_{D-1}=0\).
Thus both bottom relations hold as literal equations; they have
not been discarded by truncating the free resolution early.
The universal presentation gives a map
\(\mathcal N(X)\to P_I C\) agreeing with the constant
old map on \(\mathcal N(B)\). The pushout (136.RC4)
then gives precisely (136.RC7). This universal-property step
also shows that arbitrary existing relations in \(A\) create
no additional condition once its constant path and the stipulated
coherent boundary data have been retained. The same reasoning
works independently in each of finitely many targets.
\end{proof}

\paragraph{What this settles, and what it leaves.}
For the first possible boundary-constrained two-simplex with an
odd form-one input, the \(p=2\) native degree is \(D=4\),
and (136.RC6) permits one coefficient obstruction. For every
\(D\geq6\), that same two-simplex attachment admits the
relative path for arbitrary prescribed coherent endpoint maps,
provided its complete boundary presentation is already present.
In higher simplex dimension the exceptional dimension moves to
\(n=D-2\); it does not disappear uniformly in all simplices.

The class in (136.RC11) is asserted independent of the lift of
\(\beta\), but it is not asserted independent of the chosen
native interpolation. In the exceptional dimension its vanishing
for one interpolation would suffice. No example with a nonzero
class is claimed here. Nor is an augmentation endpoint assumed:
both endpoint maps must already extend the same map on \(A\).
These distinctions prevent the lemma from supplying, by itself,
a retained augmentation path during unrestricted horizontal
completion or an admissible integral common source.

\paragraph{Verification status.}
The author completed a whole saved-proof reading of the actual
relative strict PD pushout, its degreewise independent-block
description, the full free-resolution cone and both bottom relations,
the uniform restriction lifts, the necessary-and-sufficient correction,
and the exact relative prism degree. Root, cote1348 and cote15 completed
whole saved-proof PASS readings, recorded separately in the check register.
